\chapter{Functions}

A \emph{function} for two sets \(X\) and \(Y\) is defined as a subset of \(X \times Y\), and it is denoted by 

\[
	f: X \to Y \subseteq X \times Y 
\]

where 

\[
	\forall x \in X, !\exists y \in Y: \, (x, y) \in f.
\]

That is, for \textbf{each} \(x \in X\), there exists at most one \(y \in Y\) such that \((x, y) \in f\).
This is informally known as the vertical line test.

When writing proofs, it is important to take into consideration that the whole set \(X\) is in use by the function. This guarantees 
the existence of \(f(x)\).

\section{The Set of All Functions}

For the sets \(X\) and \(Y\), we define the set of all functions \(f: X \to Y\) as 

\[
	Y^X := \{f \mid f: X \to Y\}
\]

\subsection{Equality Of Functions}

Given two functions \(f: X \to Y\) and \(g: X' \to Y'\), they are equal if \(X = X'\) and \(Y = Y'\).

\subsection{The Characteristic Function}

Given \(T \subseteq X\), we define the \emph{characteristic function} as 

\[
	\mathcal{X}_T: X \to \{0,1\} : x \to \begin{cases} 
		1, &x \in T \\ 
		0, &x \notin T
	\end{cases}
\]

This function tells us if an element is in the subset. Also, each function \(\alpha : X \to \{0, 1\}\) defines a 
unique subset \(T\) of \(X\), 

\[
	T := \{x \in X \mid \alpha(x) = 1\}
\] 

Because of that fact, we can use this function and the concept of a subset interchangeably.

\textbf{Proof:}

We want to show \(\alpha = \mathcal{X}_T\). Therefore, we need to show that both have the same domain and image. 

For the domain, \(X \subseteq X\), \(\forall x \in X: x \in X\), which is trivial.

Now we need to prove that the images are the same; let \(x \in T \land x \in X\), then \(\alpha(x) = 1\). 
Also, if \(x \in X \setminus T\), then \(\alpha(x) = 0\). This is the definition of the characteristic function.

\QED

\section{The Graph}

We can also consider the graph as a set of ordered pairs \(x, y\) where \(y = f(x)\).

\[
	\mathrm{graph}(f):= \{(x, f(x)) \mid x \in X, f(x) \in Y\} = G
\]

\subsection{The Inverse Graph}

Let \(G\) be a graph. We define 

\[
    G^{-1} := \{(x,y) \mid (y, x) \in G\}
\]

as \emph{inverse graph} of \(G\).

\subsection{The Composite Graph}

Let \(G\) and \(H\) be graphs. We define 

\[
    G \circ H := \{(x,y) \mid \exists z \in (x, z) \in H \land (z, y) \in G\}
\]

as the \emph{composite graph} of \(G\) and \(H\).

\begin{itemize}

    \item \(G \circ (H \circ J) = (G \circ H) \circ J\).

    \item \((G^{-1})^{-1} = G\).

    \item \((G \circ H)^{-1} = H^{-1} \circ G^{-1}\).

\end{itemize}

\textbf{Proof I:}

\(((G \circ H) \circ J) = G \circ (H \circ J)\)

\[
	G \circ H = \{(z, y) \mid \exists w \in (z, w) \in H \land (w, y) \in G\} 
\]

Proceeding 

\[
	(G \circ H) \circ J = \{(x, y) \mid \exists z,w \text{ s.t. }  (x, z) \in J \land  (z, w) \in H \land (w, y) \in G  \}.
\]

For \(G \circ (H \circ J)\)

\[
	H \circ J = \{(x, w) \mid \exists z \in (x, z) \in J \land (z, w) \in H\}.
\]

Proceeding

\[
	G \circ (H \circ J) = \{(x, y) \mid \exists z,w \in  (x, z) \in J \land  (z, w) \in H \land (w, y) \in G  \} 
\]

\QED

\textbf{Proof II:}

\((G^{-1})^{-1} = G\)

\begin{align*}
	(x, y) \in (G^{-1})^{-1} &\iff (y, x) \in G^{-1} \\ 
	&\iff (x, y) \in (G^{-1})^{-1}
\end{align*}

\QED

\textbf{Proof III:}

\((G \circ H)^{-1} = H^{-1} \circ G^{-1}\).

\begin{align*}	
    (G \circ H)^{-1} &=  \{(x, y) \mid \exists z \in (y, z) \in H \land (z, x) \in G\} \\ 
	&\iff (y, z) \in H \land (z, x) \in G \\
	&\iff (z, y) \in H^{-1} \land (x, z) \in G^{-1} 
\end{align*}

\QED

\subsection{Functional Graph}

A functional graph is a graph that satisfies the second constraint of functions. Therefore, it corresponds to a function. Hence, 
all of our graph properties are applicable to functions.

\section{Domain}

The \emph{domain} of a relation \emph{R} from set \(A\) to set \(B\) is the set of all elements in 
\(A\) that are related to at least one element in \(B\). Formally, if \(R \subseteq A \times B\) 
is a relation, then the domain of \emph{R} is defined as:

\[
	\dom(R) = \{a \in A \mid \exists b \in B \text{ such that } (a,b) \in R\}
\]

In other words, the domain consists of all the input values that appear in the ordered pairs of the 
relation. For example, if \(R = \{(1,4), (2,5), (3,4), (2,6)\}\), then \(\text{Dom}(R) = \{1, 2, 3\}\).

\section{The Identity}

\[
	\text{id}(f):= idx:= \{(x, x)\mid x \in X\}
\]

\section{Image}

The \emph{image} of a relation \emph{R} from set \(A\) to set \(B\) is the set 
of all elements in \(B\) that are related to at least one element in \(A\). Formally, if 
\(R \subseteq A \times B\) is a relation, then the image of \emph{R} is defined as:

\[
	\img(R[A]) = \{y \in \text{ ran}(R) \mid \exists a \in A \text{ such that } xRy\} \subseteq B
\]

In other words, the image consists of all the output values that appear in the ordered pairs of the 
relation. For example, if \(R = \{(1,4), (2,5), (3,4), (2,6)\}\), then \(\img(R) = \{4, 5, 6\}\).

\section{The Range}

The \emph{range} of an function is the goal-set of the function. Note that this is \textbf{not} the image of the function.

\subsection{Proof That the Domain and Image Are Sets}

Given a graph \(G\) which is a set. Prove that the domain and image are also sets.

\[
	\forall (x,y) \in G: (x,y) = \{x, \{x, y\}\}
\]

Thus, the union of all the elements in \(G\) has a structure of the form

\[
	\bigcup G = \{x, y, z, \dots, \{x, y\}, \{x, z\}, \{y, z\}, \dots\}
\]

By iterating one more time we get the union of the union

\[
	\bigcup \bigcup G = \{x, y, z, \dots\}
\]

Hence, by the definition of the domain and the image of the graph we can directly see that they are subsets of the union of the union of the graph 
\(G\), and thus they are sets.

\QED

\section{The Fiber}

We denote a \emph{fiber} for \(y \in Y\) as 

\[
	f^{-1}(\{y\}) := \{x \in X \mid f(x) = y\}.
\]

\section{The Pre-Image}

Given a map \(f: X \to Y\) and a subset \(B \subseteq Y\), we denote the set 

\[
	f^{-1}(B) := \{x \in X \mid f(x) \in B\} = \bigcup_{y \in B} f^{-1}(\{y\})
\]

as the \emph{pre-image} of \(B\).

\section{Composition of Maps}

The composition of two functions is the operation of applying one function to the result of another 
function. If we have functions \(f: X \to Y\) and \(g: Y \rightarrow Z\), then the composition of \(g\) 
and \(f\), denoted as \(g \circ f\) (read as g composed with f), is a function from \(X\) to \(Z\) defined 
by:

\[
	(g \circ f)(x) = g(f(x)) \quad \text{for all } x \in X
\]

The composition applies \(f\) first, then applies \(g\) to the result. 
Note that the co-domain of \(f\) must match the domain of \(g\) for the composition to be defined.

The formal definition is 

\[
	g \circ f := \{ (x, z) \in X \times Z \mid \exists y \in Y \text{ such that } (x, y) \in f \land (y, z) \in g \}
\]

\textbf{Proof:}

To show that this is a map, take any \(x \in X\). Since \(f\) is a map from \(X\) to \(Y\), there is a unique \(y \in Y\) such that \((x,y) \in f\). Since \(g\) is a map from \(Y\) to \(Z\), there is a unique \(z \in Z\) such that \((y,z) \in g\). Thus, \((x,z) \in g \circ f\).

For uniqueness, suppose \((x,z_1), (x,z_2) \in g \circ f\). Then there are \(y_1,y_2 \in Y\) such that \((x,y_1), (x,y_2) \in f\) and \((y_1,z_1), (y_2,z_2) \in g\). Since \(f\) is a map, \(y_1 = y_2\), and since \(g\) is a map, \(z_1 = z_2\).

\QED

\textbf{Example:}

Let \(f: \R \to \R\) be defined by \(f(x) = x^2\) and \(g: \R \rightarrow \R\) be defined 
by \(g(x) = x+3\). Then:

\[
	(g \circ f)(x) = g(f(x)) = g(x^2) = x^2 + 3
\]

\[
	(f \circ g)(x) = f(g(x)) = f(x+3) = {(x+3)}^2 = x^2 + 6x + 9
\]

Note that \(g \circ f \neq f \circ g\) in general, which shows that function composition is not 
commutative.

\subsection{Associativity of the Composition}

Given the maps \(f: X \to Y\), \(g: Y \to Z\), and \(h: Z \to W\), the composition of functions is associative, i.e.,

\[
	h \circ (g \circ f) = (h \circ g) \circ f
\]

under the conditions \(f(X) \subseteq Y\) and \(g(Y) \subseteq Z\).

\textbf{Proof:}

For any \(x \in X\), we know that \(f(x) \in Y\), hence by the map property, \(\exists z \in Z, \ (f(x), z) \in g\). Continuing, 
because \(z = g(f(x)) \in Z\), there exists \(w \in W\) such that \((z, w) \in h\), and hence \(h(g(f(x))) = w\) and we have a proper map. 

Therefore, \((x, w) \in (h \circ g) \circ f\) and \((x, w) \in h \circ (g \circ f)\), proving the associativity.

\QED

\subsection{Injectivity and Surjectivity Under Composition}

Let \(f: X \to Y\) and \(g: Y \to Z\). Then the following holds: 

\begin{itemize}

	\item Is \(g \circ f\) injective, then \(f\) is injective.

	\item Is \(g \circ f\) surjective, then \(g\) is surjective.

\end{itemize}

\textbf{Proof I:}

\(g \circ f\) is injective; therefore, for \(x_1, x_2 \in X\), 

\[ 
	g(f(x_1)) = g(f(x_2)) \implies f(x_1) = f(x_2) \implies x_1 = x_2
\]

\QED

\textbf{Proof II:}

\(g \circ f\) is surjective; therefore, for all \(z \in Z\), there exists \(x \in X\) such that \(g(f(x)) = z\). This also 
implies that there exists \(y = f(x) \in Y\) such that \(g(y) = z\). 

\QED

\subsection{Domain and Range Under Composition Corollary}

Let \(G\) and \(H\) be graphs, then \(\img(H) \subseteq \dom(G) \implies \dom(G \circ H) = \dom(H)\)

\textbf{Proof:}

Note that for 

\[
	\forall y \in \img(H) : y \in \dom (G).
\]

Also, when constructing \(G \circ H\) we get a map from \(\dom(H)\) to the \(\img(G)\), because for all 
of our \(y \in \img(H)\) we can always find a pair \((y, w) \in \img(G)\) due to the structure of our graph. Hence, 
the domain of our composite graph is \(\dom(H)\).

\QED

\section{Types of Functions}

\subsection{Injective Functions}

An injective function (also called a one-to-one function) is a 
function that maps distinct elements from the domain to distinct elements in the co-domain.

Formally, a function \(f: X \to Y\) is injective if for all \(x_1, x_2 \in X\):

\[
	x_1 \neq x_2 \to f(x_1) \neq f(x_2)
\]

Equivalently, using the contrapositive:

\[
	f(x_1) = f(x_2) \to x_1 = x_2
\]

\textbf{Example:}

Prove that \(\mathrm{Id}_X: X \to X\) is an injective function.

\textbf{Proof:}

Let \(x_1, x_2 \in X\) such that \(\mathrm{Id}_X(x_1) = \mathrm{Id}_X(x_2)\). Then,

\[
	x_1 = \mathrm{Id}_X(x_1) = \mathrm{Id}_X(x_2) = x_2
\]

Hence, \(\mathrm{Id}_X\) is injective.

\QED

\subsection{Injection}

Given \(X \subseteq Y\), the \emph{injection} from \(X\) to \(Y\) defined as 

\[
	\iota_{X}: X \to Y : x \mapsto x
\]

is an \emph{injective function}.

\subsection{Surjective Functions}

A \emph{surjective function} (also called an onto function) is a function whose image equals its co-domain, 
meaning that every element in the co-domain has at least one pre-image in the domain.

Formally, a function \(f: X \to Y\) is surjective if:

\[
	\forall y \in Y, \exists x \in X \text{ such that } f(x) = y
\]

To prove that a function is surjective, one must show that given any element \(y \in Y\), there exists an element \(x \in X\) such that \(f(x) = y\). When 
a concrete function definiton like \(f(x) = 2x\) is given we just need to solve for \(x\), and show that it is an element of the domain. For the case 
where we are dealing with an more abstract function, we need to use the \textbf{function definition} to guarantee the existence of elements in the domain.

\textbf{Example I:}

Prove that the function \(f: \R \to \R\) defined by \(f(x) = 2x\) is surjective.

\textbf{Proof:}

Let \(y \in \R\), so \(y = 2x\). Solving for \(x\) gives \(x = \frac{y}{2}\). Note that \(\frac{y}{2} \in \R\), hence 
we have shown that for any \(y \in \R\), there exists an \(x \in \R\) such that \(f(x) = y\). Therefore, \(f\) is surjective.

\QED

\textbf{Example II:}

Given the characteristic function \(\mathcal{X}_T: X \to \{0,1\}\) defined by

\[
	\mathcal{X}_T(x) = 
	\begin{cases}
		1 & \text{if } x \in T \\
		0 & \text{if } x \notin T
	\end{cases}
\]

where \(T \subseteq X\).

Prove that \(\mathcal{X}_T\) is surjective. 

\textbf{Proof:}

Let \(b \in \{0,1\}\). We need to show that there exists an \(x \in X\) such that \(\mathcal{X}_T(x) = b\).

If \(b = 1\), then choose \(x \in T\). By definition, \(\mathcal{X}_T(x) = 1\).

If \(b = 0\), then choose \(x \in X \setminus T\). By definition, \(\mathcal{X}_T(x) = 0\).

Hence, for every \(b \in \{0,1\}\), there exists an \(x \in X\) such that \(\mathcal{X}_T(x) = b\). Therefore, \(\mathcal{X}_T\) is surjective.

\QED

\subsection{Bijective Functions}

A \emph{bijective function} (also called a one-to-one correspondence) is a function that is both injective and surjective. In other words, every 
element in the co-domain is mapped to by exactly one element in the domain.

Formally, a function \(f: X \to Y\) is bijective if it is both:

\begin{itemize}

	\item \emph{Injective}: \(\forall x_1, x_2 \in X, x_1 \neq x_2 \to f(x_1) \neq f(x_2)\)

	\item \emph{Surjective}: \(\forall y \in Y, \exists x \in X \text{ such that } f(x) = y\)

\end{itemize}

Bijective functions establish a perfect pairing between elements of the domain and co-domain, where each 
element in the domain corresponds to exactly one element in the co-domain, and vice versa. A bijection 
allows us to define an inverse function \(f^{-1}: Y \to X\).

\section{Cardinality of Injective and Surjective Functions}

Given finite sets \(X\) and \(Y\) with the map \(f: X \to Y\), the following holds: 

\begin{itemize}

	\item If \(f\) is injective, then \(|X| \le |Y|\). 

	\item If \(f\) is surjective, then \(|X| \ge |Y|\).

	\item If \(f\) is bijective, then \(|X| = |Y|\).

\end{itemize}

\textbf{Proof I:}

\(f\) is injective; therefore, \(|f^{-1}(y)| \le 1\) for all \(y \in Y\), hence,

\[
	|X| = \sum_{y \in Y} |f^{-1}(y)| \le \sum_{y \in Y} 1 = |Y|.
\]

\QED

\textbf{Proof II:}

\(f\) is surjective, so \(|f^{-1}(y)| \ge 1\) for all \(y \in Y\); therefore,

\[
	|X| = \sum_{y \in Y} |f^{-1}(y)| \ge \sum_{y \in Y} 1 = |Y|.
\]

\QED

\textbf{Proof III:}

Proven by using I and II. 

\QED

\section{Propositions on Images and Pre-Images Under Set Operations}

Let \( f : X \to Y \) be a function.

\subsection{Union and Cut Sets}

\begin{itemize}

	\item For subsets \( A_1, A_2 \subseteq X \), we have:

	      \[
		      f(A_1 \cup A_2) = f(A_1) \cup f(A_2)
		      \quad \text{and} \quad
		      f(A_1 \cap A_2) \subseteq f(A_1) \cap f(A_2)
	      \]

	\item For subsets \( B_1, B_2 \subseteq Y \), we have:
	      
	      \[
		      f^{-1}(B_1 \cup B_2) = f^{-1}(B_1) \cup f^{-1}(B_2)
		      \quad \text{and} \quad
		      f^{-1}(B_1 \cap B_2) = f^{-1}(B_1) \cap f^{-1}(B_2)
	      \]
\end{itemize}

\textbf{Proof:}

We prove the second part of 1 and the first part of 2.

Let \( y \in f(A_1 \cap A_2) \). Then there exists \( x \in A_1 \cap A_2 \) such that 
\( f(x) = y \). Since \( x \in A_1 \) and \( x \in A_2 \), it follows that \( y \in f(A_1) \) 
and \( y \in f(A_2) \), hence \( y \in f(A_1) \cap f(A_2) \). Therefore, every element of 
\( f(A_1 \cap A_2) \) is also an element of \( f(A_1) \cap f(A_2) \), so:

\[
	f(A_1 \cap A_2) \subseteq f(A_1) \cap f(A_2)
\]

The reason we do not have equality is because: 

\begin{align*}
	y \in f(A_1) \cap f(A_2) &\iff \exists f^{-1}(y) \in A_1 \land f^{-1}(y) \in A_2 \\
	&= x_1 \in A_1 \land x_2 \in A_2
\end{align*}

Note, that if \(x_1 = x_2\) we have equality, but if \(x_1 \neq x_2\) then \(y \notin f(A_1 \cap A_2)\). Hence 
we have a propert subset.

\QED

Let \( x \in f^{-1} (B_1 \cup B_2) \). Then \( f(x) \in B_1 \cup B_2 \), which means \( f(x) \in B_1 \) 
or \( f(x) \in B_2 \). Thus, \( x \in f^{-1}(B_1) \) 
or \( x \in f^{-1}(B_2) \), which implies:
	      
\[
	 x \in f^{-1}(B_1) \cup f^{-1}(B_2)
\]

Hence, both sets contain the same elements and are therefore equal:
	      
\[
	f^{-1}(B_1 \cup B_2) = f^{-1}(B_1) \cup f^{-1}(B_2)
\]

\QED

\subsection{Union and Cut of the Whole Domain and Range}

Let \( f : X \to Y \) be a function.

\begin{itemize}

	\item Let \( \mathcal{F} \) be a collection of subsets of \(X\). Then:

   	      \[
		      f\left( \bigcup_{A \in \mathcal{F}} A \right) = \bigcup_{A \in \mathcal{F}} f(A)
		      \quad \text{and} \quad
		      f\left( \bigcap_{A \in \mathcal{F}} A \right) \subseteq \bigcap_{A \in \mathcal{F}} f(A)
	      \]

	\item Let \( \mathcal{G} \) be a collection of subsets of \emph{Y}. Then:

	      \[
		      f^{-1}\left( \bigcup_{B \in \mathcal{G}} B \right) = \bigcup_{B \in \mathcal{G}} f^{-1}(B)
		      \quad \text{and} \quad
		      f^{-1}\left( \bigcap_{B \in \mathcal{G}} B \right) = \bigcap_{B \in \mathcal{G}} f^{-1}(B)
	      \]

\end{itemize}

\textbf{Proof (partial):}

We show the first statement of part (ii); the rest follows analogously.

Let \( x \in f^{-1} \left( \bigcup_{B \in \mathcal{G}} B \right) \). Then:

\[
	f(x) \in \bigcup_{B \in \mathcal{G}} B
	\quad \Leftrightarrow \quad
	\exists B \in \mathcal{G} \text{ such that } f(x) \in B
	\quad \Leftrightarrow \quad
	\exists B \in \mathcal{G} \text{ such that } x \in f^{-1}(B)
\]

Hence:

\[
	x \in \bigcup_{B \in \mathcal{G}} f^{-1}(B)
\]

It follows that:

\[
	f^{-1} \left( \bigcup_{B \in \mathcal{G}} B \right) = \bigcup_{B \in \mathcal{G}} f^{-1}(B)
\]

\QED

\textbf{Example:}

We define \(A = \{1, 0\}\) and \(B = \{2\}\) and \(f(1) = 3\), \(f(2) = 3\), \(f(0) = 1\). Then 

\[
	f(A \cap B) = f( \emptyset) = \emptyset
\]

And 

\[
	f(A) \cap f(B) = 3
\]

\section{Bijection of the Power Set and \texorpdfstring{\(2^A\)}{}}

We denote \(2^A\) as the set of all functions from \(A\) to a set of two elements.
The statement goes as follows: there is a bijection between the power set \(\mathcal{P}(A)\) and 
\(2^A\).

\textbf{Proof:}

Let \(\gamma: \mathcal{P}(A) \to 2^A\). If \(B \in \mathcal{P}\) let 

\[
	\gamma (B) = C_B, \forall B \in \mathcal{P}(A).
\]

By this definition we get the properties of a function directly. Now we have to prove that 
the function is a one-to-one correspondence.

For the injectivity, we already know that for \(\gamma(B) = \gamma(C) \implies B = C\), since otherwise it would be the case 
that we have two elements in the power set which are identical.

For the surjectivity, if \(f \in 2^A\), and we let \(B = f^{-1}(0)\), then \(f = C_B = \gamma(B)\). Thus, 
\(\gamma\) satisfies the surjectivity.

\QED

\section{Restriction Function}

Let \(G\) be a graph and \(B \subseteq \dom(G)\) then we denote the \emph{restriction} of \(G\) as:

\[
	G_{[B]} = \{(x, y) \mid (x, y) \in G \land x \in B\}.
\]

The range of this restricted map can be denoted as:

\[
	G(B) = \{y \mid x \in \dom(B) \land )(x, y) \in G\}
\]

\section{Even Function}

A real function \(f\) is \emph{even} if

\begin{itemize}
	
	\item \(\forall x \in \text{Dom}(f) \exists -x \in \text{Dom}(f)\)
	
	\item \(\forall x \in \text{Dom}(f) \text{ the following is true } f(x) = f(-x)\)

\end{itemize}

\section{Odd Function}

A real function \(f\) is \emph{odd} if

\begin{itemize}
	
	\item \(\forall x \in \text{Dom}(f) \exists -x \in \text{Dom}(f)\)
	
	\item \(\forall x \in \text{Dom}(f) \text{ the following is true } f(x) = - f(x)\)

\end{itemize}

\section{Inverse of a Function}

The \emph{inverse} of a function \(f: X \to Y\) is a function \(f^{-1}: Y \to X\) that 
reverses the operation of \(f\). That is, if \(f\) maps an element \(x \in X\) to an 
element \(y \in Y\), then the inverse function \(f^{-1}\) maps \(y\) back to \(x\).

Formally, a function \(f: X \to Y\) has an inverse \(f^{-1}: Y \to X\)  
if and only if \(f\) is bijective (both injective and surjective). The inverse function satisfies the 
following properties:

\[
	f^{-1}(f(x)) = x \quad \text{for all } x \in X,
\]

\[
	f(f^{-1}(y)) = y \quad \text{for all } y \in Y.
\]

In other words, composing a function with its inverse yields the identity function. That is:

\[
	f^{-1} \circ f = id_X \quad \text{and} \quad f \circ f^{-1} = id_Y
\]

where \(id_X\) and \(id_Y\) are the identity functions on sets \(X\) and \(Y\), respectively.

\textbf{Proof:}

\(\Rightarrow\): direction is trivial, due to the definition of a bijective function. 

\(\Leftarrow\): Our function has the following properties:

\[
	f^{-1}(f(x)) = x \quad \text{for all } x \in X,
\]

\[
	f(f^{-1}(y)) = y \quad \text{for all } y \in Y.
\]

The first condition gives us injectivity because for all \(y \in Y\) there exists exactly one \(x \in X\)
such that \(f(x) = y\). Otherwise, the condition would not be obeyed.

The second one gives us the surjectivity of \(f\) by providing a definition for each \(y \in Y\). Otherwise, 
our condition would be violated. 

The uniqueness is guaranteed by 

\[
	f^{-1}(y) \ne g^{-1}(y)
\]

This gives us a contradiction if they were to have the same properties but were different.

\QED

\subsection{Properties of the Inverse Function}

Given \(f: X \to Y\), with \(A \subseteq X\) and \(B \subseteq Y\). The following properties hold for the inverse function \(f^{-1}\):

\begin{itemize}

	\item \(f^{-1}(f(A)) = A\).

	\item \(f(f^{-1}(B)) = B\).

	\item \(f(f^{-1}(B)) \subseteq B\).

	\item \(A \subseteq f^{-1}(f(A))\).

\end{itemize}

We will prove the last two statements, since the fist two are trivial due to the function definition.

\textbf{Proof III:}

By definition 

\[
	f^{-1}(B) = \{x \in X \mid f(x) \in B\}
\]

Note that there may be some \(b \in B\) such that \(b \neq f(x) \forall x \in f^{-1}(B)\), thus proving that 
\(f(f^{-1}(B)) \subseteq B\).

\QED

\textbf{Proof IV:}

By definition 

\[
	f(A) = \{y \in Y \mid f^{-1}(x) \in A\}
\]

Similarly, there might be some \(y \in f(A)\) for which there exits some \(x \in X \setminus A\) such that \(f(x) \in f(A)\). Hence, 
we do not have a set equality.

\QED

\subsection{Steps to Find the Inverse of a Function}

To find the inverse of a function \(f (x)\), follow these steps:

\begin{enumerate}
	
	\item Replace \(f (x)\) with \(y\): \(y = f (x)\)
	
	\item Interchange the variables \(x\) and \(y\): \(x = f (y)\)
	
	\item Solve for \(y\) in terms of \(x\): \(y = f^{-1} (x)\)
	
	\item Verify that the resulting function is indeed the inverse by checking that \(f^{-1}(f(x)) = x\) 
		  and \(f (f^{-1} (x)) = x\)

\end{enumerate}

\textbf{Example:}

Let's apply the steps above to find the inverse of \(f (x) = 2x + 3\).

\textbf{Step 1: Replace \(f(x)\) with \(y\)}

\[
	y = 2x + 3
\]

\textbf{Step 2: Interchange the variables \(x\) and \(y\)}

\[
	x = 2y + 3
\]

\textbf{Step 3: Solve for \(y\) in terms of \(x\)}

\begin{align*}
	x               & = 2y + 3 \\
	x - 3           & = 2y     \\
	\frac{x - 3}{2} & = y
\end{align*}

So, the inverse function is:

\[
	f^{-1}(x) = \frac{x - 3}{2}
\]

\textbf{Step 4: Verify that \(f^{-1} (f (x)) = x\) and \(f(f^{-1} (x)) = x\)}

Let's verify \(f^{-1} (f (x)) = x\):

\begin{align*}
	f^{-1}(f(x)) & = f^{-1}(2x + 3)         \\
	             & = \frac{(2x + 3) - 3}{2} \\
	             & = \frac{2x}{2}           \\
	             & = x
\end{align*}

And let's verify \(f(f^{-1}(x)) = x\):

\begin{align*}
	f(f^{-1}(x)) & = f\left(\frac{x - 3}{2}\right)     \\
	             & = 2\left(\frac{x - 3}{2}\right) + 3 \\
	             & = (x - 3) + 3                       \\
	             & = x
\end{align*}

Since both compositions yield the identity function, \(f^{-1} (x) = \frac{x - 3}{2}\) is indeed the 
inverse of \(f (x) = 2x + 3\).

\subsection{Properties of the Inverse Function}

Let \(f: X \to Y\) be a function.

\begin{itemize}

	\item Assume \(x \sim y\) if \(f(x) = f(y)\); then \(\sim\) is an equivalence relation.

	\item Consider the quotient set \( X_f := X/\sim \), where \( \sim \) is the equivalence relation 
	      defined by \( x \sim x' \iff f(x) = f(x') \). Let \( q_f : X \to X_f \) be the canonical 
		  projection defined by \( q_f(x) = [x]_{\sim} \), and let \( \iota_f : f(X) \to Y \) be the inclusion  
		  map, \( y \mapsto y \). Then the function
	      
		  \[
		      \hat{f} : X_f \to f(X), \quad [x]_{\sim} \mapsto f(x)
	      \]
	  
		  is a bijection, and the original map \(f\) can be written as the composition
	 
		  \[
		      f = \iota_f \circ \hat{f} \circ q_f.
	      \]
	 
		  Here, \( q_f \) is surjective, \( \hat{f} \) is bijective, and \( \iota_f \) is injective. 
		  This yields the following commutative diagram:
	      
		  \begin{center}
				\begin{tikzcd}
					X \arrow[r, "f"] \arrow[d, "q_f"] & Y                         \\
					X_f \arrow[r, "\hat{f}"]          & f(X) \arrow[u, "\iota f"]
				\end{tikzcd}
	      \end{center}

\end{itemize}

\subsection{Left and Right Inverses}

Given \(f: X \to Y\), \(g: Y \to X\), and \(\gamma: Y \to X\), then  

\begin{itemize}
	
	\item \(f\) is \emph{injective} if and only if \(g \circ f = \mathrm{Id}_X\). Each \(g\) is surjective and called the \emph{left inverse} of \(f\).


	\item \(f\) is \emph{surjective} if and only if \(f \circ \gamma = \mathrm{Id}_Y\). Each \(\gamma\) is injective and called the \emph{right inverse} of \(f\).
	
\end{itemize}

\textbf{Proof I:}

\(\Leftarrow\): We assume \(g \circ f = \mathrm{Id}_X\). Let \(f(x_1) = f(x_2)\) for \(x_1, x_2 \in X\). Applying our composite function gives 
\(g(f(x_1)) = g(f(x_2)) \implies x_1 = x_2\), since \(g \circ f = \mathrm{Id}_X\). Therefore, \(f\) is injective.

\(\Rightarrow\): Let \(f\) be injective; now let us define 

\[
	g: 
	\begin{cases}
		x, & f(x) = y \\ 
		x_0, & y \notin \img(f) 
	\end{cases}
\]

For \(x \in \dom(f)\), \(g \circ f = \mathrm{Id}_X\) and \(g\) is also surjective since for all \(y \in Y\) there exists a pre-image which is not empty.

\QED

\textbf{Proof II:}

\(\Rightarrow\): Assume \(f \circ \gamma = \mathrm{Id}_Y\); now let \(y \in Y\), then \(f(\gamma(y)) = y \implies y \in \img(f)\), so \(f\) is 
surjective.

\(\Leftarrow\): Let \(f\) be surjective. This means that \(|f^{-1}(\{y\})| \neq 0\) for all \(y \in Y\). Now let us define 

\[
	\gamma:
	\begin{cases}
		x, & |f^{-1}(\{y\})| = 1 \\
		\mathrm{choose}(f^{-1}(\{y\})), & \text{ else }
	\end{cases}
\]


Note that \(f \circ \gamma = \mathrm{Id}_Y\) since there is always at least one \(x\) such that we return to \(Y\), and if there are multiple choices, 
we choose one from the fiber.

\QED

\textbf{Example:}

Given \(f: X \to Y\) having a right inverse \(\gamma: Y \to X\) guarantees the existence of a solution to \(f(x) = y\). Set \(x = \gamma(y)\), then 
\(f(x) = f(\gamma(y)) = (f \circ \gamma)(y) = \text{Id}_Y(y) = y\).

\subsection{Equivalent Definitions of the Inverse Function}

Given the map \(f: X \to Y\), both statements are equivalent 

\begin{enumerate}

	\item \(f^{-1}:= \{(y, x) \in Y \times X \mid (x, y) \in f\}\) is a map from \(Y \to X\).

	\item  \(\forall y \in Y, \, |f^{-1}(y)| = 1\).

\end{enumerate}

\textbf{Proof:}

\(\Rightarrow\): We assume that 1. is true, hence \(\forall y \in Y, !\exists x \in X: f^{-1}(y) = x\), in other words 

\[
	f^{-1}(\{y\}) = \{x\}.
\]

Thus, \(f^{-1}(\{y\}) = \{x\}\), and the claim follows.

\(\Leftarrow\): For the other direction, if the fiber only has one element, then \(f^{-1}\) is a map.

\QED

\subsection{Inverse Definition over Left and Right Inverses}

A function is \emph{bijective} (invertible) is and only if it has a left and right inverse.

\textbf{Proof:}:

\(\Rightarrow\): Our function is bijective, and we know that by the equivalences that it has to have both a left and a right inverse.

\(\Leftarrow\): If a function has both a left and a right inverse, it must be bijective. The existence of a left inverse ensures injectivity, and the existence of a right inverse ensures surjectivity. Therefore, the function is bijective.

\QED

\section{Transformations of a Function}

\emph{Transformations} modify the appearance of a function's graph without altering its basic shape.
Here, we examine how different algebraic changes to a function \( f(x) \) affect its graph:

\emph{Vertical Translation:} \( f(x) + a \)
	    
\begin{itemize}
	
	\item Shifts the graph \emph{upward} if \( a > 0 \), and \emph{downward} if \( a < 0 \).
	
	\item Each point on the graph moves vertically by \(a\) units.

\end{itemize}

\emph{Horizontal Translation:} \( f(x + a) \)

\begin{itemize}

	\item Shifts the graph \emph{left} if \( a > 0 \), and \emph{right} if \( a < 0 \).

	\item This is opposite of what might be expected: adding to \(x\) shifts the graph in the negative 
	      direction.

\end{itemize}

\emph{Vertical Scaling (Stretch/Compression):} \( a f(x) \)

\begin{itemize}
	
	\item If \( |a| > 1 \): the graph is \emph{stretched} vertically (taller and narrower).
	
	\item If \( 0 < |a| < 1 \): the graph is \emph{compressed} vertically (shorter and wider).
	
	\item If \( a < 0 \): includes a reflection across the \emph{x-axis}.

\end{itemize}

\emph{Horizontal Scaling (Stretch/Compression):} \( f(a x) \)

\begin{itemize}

	\item If \( |a| > 1 \): the graph is \emph{compressed} horizontally (narrower).

	\item If \( 0 < |a| < 1 \): the graph is \emph{stretched} horizontally (wider).

	\item If \( a < 0 \): includes a reflection across the \emph{y-axis}.

\end{itemize}

\emph{Reflection across the x-axis:} \( -f(x) \)
	
\begin{itemize}

	\item Flips the graph upside-down over the x-axis.

	\item Each point \( (x, y) \) becomes \( (x, -y) \).

\end{itemize}

\emph{Reflection across the y-axis:} \( f(-x) \)

\begin{itemize}

	\item Flips the graph left-to-right over the y-axis.

	\item Each point \( (x, y) \) becomes \( (-x, y) \).

\end{itemize}

\section{Counting Functions From \texorpdfstring{\(X \to Y\)}{X to Y}}

Given \(f: X \to Y\) with \(|X| = n\) and \(|Y| = m\), we can count the number of different types of functions from 
\(X\) to \(Y\):

\begin{itemize}

	\item \emph{Total functions:} There are \(m^n\) total functions from \(X\) to \(Y\). Each of the \(n\) elements in 
		  \(X\) can be mapped to any of the \(m\) elements in \(Y\).

	\item \emph{Injective functions:} If \(n \leq m\), the number of injective functions is given by 
	\(m \times (m-1) \times (m-2) \times \ldots \times (m-n+1)\). This is because the first element in \(X\) has \(m\) choices, the second has \(m-1\)
	choices, and so on.

	\item \emph{Surjective functions:} The number of surjective functions can be calculated using the 
		  principle of inclusion-exclusion:
		  
		  \[
			  \text{Number of surjective functions} = \sum_{k=0}^{m} (-1)^k \binom{m}{k} (m-k)^n
		  \]
		  
		  This formula counts all functions and subtracts those that miss at least one element in \(Y\). 

	\item \emph{Bijective functions:} If \(n = m\), the number of bijective functions (one-to-one and onto) is 
		  \(n!\) (n factorial), which counts all possible permutations of the \(n\) elements in \(X\) to the \(n\) elements in \(Y\).

\end{itemize}